\usepackage[T1]{fontenc}
\usepackage{lmodern}
\usepackage{amsmath,amssymb,amsthm,mathtools,mathrsfs}
\usepackage[margin=1.05in]{geometry}
\usepackage{microtype}
\usepackage{enumitem,booktabs,graphicx}
\usepackage{tikz,tikz-cd}
\usetikzlibrary{arrows.meta,positioning,calc}
\usepackage{xcolor}
\definecolor{linkblue}{RGB}{22,53,93}
\usepackage[colorlinks=true,linkcolor=linkblue,citecolor=linkblue,urlcolor=linkblue]{hyperref}
\usepackage[capitalise,nameinlink]{cleveref}
\input{glyphtounicode}
\pdfgentounicode=1
\input{glyphtounicode-cmex}

\let\manuscriptoriginalhookrightarrow\hookrightarrow
\protected\def\hookrightarrow{\mathrel{%
  \pdfliteral direct{/Span << /ActualText <FEFF21AA> >> BDC}%
  \manuscriptoriginalhookrightarrow
  \pdfliteral direct{EMC}}}
\let\manuscriptoriginallongmapsto\longmapsto
\protected\def\longmapsto{\mathrel{%
  \pdfliteral direct{/Span << /ActualText <FEFF27FC> >> BDC}%
  \manuscriptoriginallongmapsto
  \pdfliteral direct{EMC}}}

\hypersetup{pdftitle={Cyclic length and chromatic fixed-point loss},pdfauthor={OpenAI}}
\setlength{\emergencystretch}{2em}
\setlist{itemsep=3pt,topsep=5pt}
\numberwithin{equation}{section}
\newtheorem{theorem}{Theorem}[section]
\newtheorem{lemma}[theorem]{Lemma}
\newtheorem{proposition}[theorem]{Proposition}
\newtheorem{corollary}[theorem]{Corollary}
\theoremstyle{definition}
\newtheorem{definition}[theorem]{Definition}
\newtheorem{example}[theorem]{Example}
\theoremstyle{remark}
\newtheorem{remark}[theorem]{Remark}
\newcommand{\Sp}{\mathrm{Sp}}
\DeclareMathOperator{\CAlg}{CAlg}
\DeclareMathOperator{\Mod}{Mod}
\DeclareMathOperator{\Fun}{Fun}
\DeclareMathOperator{\Map}{Map}
\DeclareMathOperator{\Hom}{Hom}
\DeclareMathOperator{\End}{End}
\DeclareMathOperator{\Aut}{Aut}
\DeclareMathOperator{\Sur}{Sur}
\DeclareMathOperator{\im}{im}
\DeclareMathOperator{\rank}{rank}
\DeclareMathOperator{\supp}{supp}
\DeclareMathOperator{\Spc}{Spc}
\DeclareMathOperator{\tr}{tr}
\DeclareMathOperator{\Bor}{Bor}
\newcommand{\id}{\mathrm{id}}
\newcommand{\Q}{\mathbb Q}
\newcommand{\Z}{\mathbb Z}
\newcommand{\N}{\mathbb N}
\newcommand{\F}{\mathbb F}
\newcommand{\one}{\mathbf1}
\newcommand{\cts}{\mathrm{cts}}
\newcommand{\op}{\mathrm{op}}
